\documentclass[11pt]{article}
\usepackage[margin=1in]{geometry}
\usepackage{amsmath,amssymb,bm}
\title{Energy flow and power exchange}
\author{}
\date{}
\begin{document}
\maketitle

\section{Problem Setup}
The exterior calculation supplies an in-plane electric field
$\bm E=(E_x,E_y,0)$ throughout the mesh. The magnetic calculation
supplies $\bm B=(0,0,B_z)$. We want to show where electromagnetic
energy flows and where the battery supplies energy that the
conductor absorbs as heat.

\section{Deriving the energy-flow vector}
In this model $\bm H=\bm B/\mu$. The Poynting vector is the
cross product
\begin{equation}
 \bm S=\bm E\times\bm H
      =\frac{B_z}{\mu}(E_y,-E_x,0).
 \label{eq:poynting}
\end{equation}
The electric field is constant on each triangle, while $B_z$
is linear there. Thus $\bm S$ is linear within a triangle.
Its exact triangle average is
\begin{equation}
 \overline{\bm S}_h
 =\frac{\overline B_{z,h}}{\mu}(E_{y,h},-E_{x,h}),
 \qquad
 \overline B_{z,h}
 =\frac{b_{h,0}+b_{h,1}+b_{h,2}}{3}.
\end{equation}
The app draws arrows from these averages. They describe energy
flow, not the motion of electrons. In the strictly two-dimensional
model, $B_z=0$ in the outside air, so $\bm S=0$ there even though
$\bm E$ may be nonzero. The cavity can have nonzero energy flow
because its magnetic field is generally nonzero.

To connect this flow to the circuit, use the vector identity
\begin{equation}
 \nabla\cdot(\bm E\times\bm H)
 =\bm H\cdot(\nabla\times\bm E)
  -\bm E\cdot(\nabla\times\bm H).
\end{equation}
In steady state $\nabla\times\bm E=0$, and Ampere's law gives
$\nabla\times\bm H=\bm J$. Therefore
\begin{equation}
 \nabla\cdot\bm S=-\bm J\cdot\bm E.
 \label{eq:balance}
\end{equation}
The relation says that energy leaving a small volume through
the field is balanced by electrical work inside it.

\section{Power supplied by the battery and absorbed as heat}
Inside the conductor, the current law is
$\rho\bm J=\bm E+\bm F$. Rearranging and taking the dot
product with $\bm J$ gives
\begin{equation}
 \bm J\cdot\bm E
 =\rho|\bm J|^2-\bm J\cdot\bm F.
 \label{eq:work}
\end{equation}
The term $\rho|\bm J|^2$ is Joule heating in the wire,
battery, and resistors. The term $\bm J\cdot\bm F$
is power supplied by the battery's driving field and
is zero outside the battery. The app's power-exchange
view shows
\begin{equation}
 p_{\rm exchange,h}
 =\bm J_h\cdot\bm F_h-\rho_h|\bm J_h|^2.
\end{equation}
Positive values mean net supply, while negative values
mean net absorption as heat. In air, current and both
power-density terms are zero.

The battery-power and resistor-heating readouts sum the
corresponding triangle densities times area. Let $\mathcal T_B$ be
the set of battery triangles and $\mathcal T_R$ the set of resistor
triangles:
\begin{equation}
 P'_{\rm battery}
 =\sum_{h\in\mathcal T_B}A_h\,\bm J_h\cdot\bm F_h,
 \qquad
 P'_{\rm resistors}
 =\sum_{h\in\mathcal T_R}A_h\,\rho_h|\bm J_h|^2.
\end{equation}
These powers are per unit depth in the unchanged $z$
direction. Since EMF is work per unit charge, the displayed
loop current is $I'=P'_{\rm battery}/\mathcal E$ when
$\mathcal E\ne0$.

\paragraph{Code:} \texttt{calculations/poynting.py}.
\end{document}
